\subsection{固定域上的学习误差与条件预算}

为将新增变质量模型的轨迹误差进一步联系到KAM所需的函数误差，在计算前固定其正则状态域$D_{q,p}=[-0.25,0.25]^4$，其覆盖全部保存训练、验证与测试状态及中点输入（最大绝对分量$0.199230$）。在$X=f(q)=q+q^3$、$Y=p/f'(q)$的坐标中，已知物理主部严格相消，能量误差等于网络残差减去$V_{\rm nl}(f(q))$。记该误差为$e_\theta$，以常数平移对齐能量，并在$16^4$复合中点网格和独立扰动Sobol点上评价
\[
 \mathcal E_6=\left(\frac1{|D_{q,p}|}\int_{D_{q,p}}
       \sum_{j=1}^4|\partial_j e_\theta|^6\,dx\right)^{1/6}.
\]
这里采用与\citet[式(8)]{Chen2022}一致的分量$6$范数，$6>4$；未归一化的Lebesgue $L^6$量需另乘$|D_{q,p}|^{1/6}=0.629961$。表\ref{tab:kam_error_bridge}的$\widehat{\mathcal E}_6$为$2^{16}$个Sobol点的域平均估计，不是训练分布上的泛化界；原训练仍采用二次损失。该量采用新$(q,p)$坐标及其域测度，不与旧$(X,Y)$盒域的数值直接横比。

为避免把有限点最大值当作全域界，进一步利用$|\tanh'|\leq1$、$|\tanh''|\leq4/(3\sqrt3)$及保存权重，逐分量界定误差Hessian。物理残差的导数在新坐标中重新计算，包含由$f''$产生的链式项。以$17^4$端点网格的最近点Taylor估计覆盖单元，并计入复合中点均值与真实域均值之间的求积误差，得到零均值能量误差的保守包络$U_{\rm learn}^{(0)}$。该包络由实数算术下的解析不等式构造，数值计算未使用向外舍入区间，故属于可复算保守包络，不构成机器区间认证。

\begin{table}[t]
\centering\small
\caption{固定域的误差诊断。前两行报告五种子均值$\pm$样本标准差，后两行为五种子最大保守包络；能量均为无量纲。}
\label{tab:kam_error_bridge}
\begin{tabular}{lrr}
\toprule
量 & GR & GC\\
\midrule
$\widehat{\mathcal E}_6\;(10^{-5})$ & $8.660\pm1.540$ & $8.121\pm1.068$\\
采样能量最大误差$(10^{-5})$ & $2.512\pm0.597$ & $1.944\pm0.371$\\
$U_{\rm learn}^{(0)}\;(10^{-4})$ & $4.047$ & $2.823$\\
$U_{\rm phys}^{(0)}+U_{\rm learn}^{(0)}\;(10^{-4})$ & $4.107$ & $2.883$\\
\bottomrule
\end{tabular}
\end{table}

二次参考$h_0$的退化性促使我们另将已有角平均定义为显式可积比较参考$K(I)=\nu^TI+\overline h_4(I)$，其中$I_j=(X_j^2+Y_j^2)/2$仍在原物理模态坐标中定义，取$I_j\in[0.0008,0.02]$，对应参考半径$0.04$--$0.20$。其Hessian恒为$B_4$，最小奇异值为$5.80643\times10^{-4}$，因此该\emph{比较参考}在所取域上非退化。完整物理扰动为$R_{\rm phys}=h_*-K$，仍保留角依赖四次项；其$C^0$保守包络为$U_{\rm phys}^{(0)}=5.98892\times10^{-6}$。单独的余弦六阶Taylor余项仅为$2.39336\times10^{-9}$，不能替代完整$R_{\rm phys}$。有限整数集$0<|k|_1\leq32$上，对该参考全作用量矩形的频率区间筛查下界为$0.152123$，该值不延伸到全整数格的Diophantine条件。该作用量环域映回新坐标后包含于$D_{q,p}$，故$C^0$界可直接搬运；高阶范数还需计入变换Jacobian及其导数。

上述量使学习误差与物理扰动的$C^0$预算可以逐项比较，但尚未满足KAM定理的全部前提。对于实际数值预测，还需在所选定理规定的解析或高阶范数$\mathcal N$下建立
\[
 B_{\mathcal N}^{\rm phys}+B_{\mathcal N}^{\rm learn}
       +B_{\mathcal N}^{\rm disc}(h)<\epsilon_*,
\]
并核验相应图册、非共振和阈值常数；表中的$C^0$包络不能直接替代这些量。二次主部的精确Cayley修改频率可用于分离步长量级，却不是完整非线性格式的修改Hamiltonian界。\citet[定理11]{Chen2022}另要求目标可积及规定的泛化条件，本研究未援用其概率环面结论。因而本节补充的是可复算的误差预算及其与条件之间的缺口，未获得严格不变环面证书。
